% ===== 原文 pp.28（末段）–30 =====
% 译者 tr-04：§4.2 前半（p30 含 4.2.1 行内小标题、式 (4.5)(4.6)）

% [原文 p.28]
仿照 [59] 与 [22]，人们可以定义关于无不动点对合 $\tau$ 的反等变 Higgs 丛（anti-equivariant Higgs bundles）$(\bar\partial,\Phi)$，即由 $\tau$-等变的全纯结构 $\bar\partial$ 与满足
\[
\tau^*\Phi=-g^{-1}\circ\Phi\circ g
\]
的 Higgs 场 $\Phi$ 组成的对，其中 $g$ 是满足 $\tau^*\bar\partial=\bar\partial\cdot g$ 与 $(g\circ\tau)g=\id$ 的规范变换。

% [原文 p.29]
这些 Higgs 丛给出 Higgs 丛模空间中拉格朗日子空间的例子；虽然它们对任意亏格的紧黎曼面都有定义，但这里我们按照 [59] 只关注奇亏格的情形（更一般的设定在 [22] 中处理）。事实上，[59, Proposition 15] 表明：$\cM_{\Sl(2,\C)}$ 中反等变 $\Sl(2,\C)$-Higgs 丛的连通分支关于复结构 $I$ 是复子流形、关于复结构 $J$ 与 $K$ 是拉格朗日的，从而是 $(\mathrm{B},\mathrm{A},\mathrm{A})$-膜。最后，我们应当提到，当考虑把 [59] 中的结果推广到更一般的群时，会出现有趣的研究方向。

% [Tutor Note: p.29 原文两处笔误：(1) "Lagrangians subspaces" 应为 "Lagrangian subspaces"；(2) "it is shown in [59, Proposition 15] that anti-equivariant SL(2,C)-Higgs bundles for connected components … which are complex submanifolds … and thus define (B,A,A)-branes" 一句缺谓语（"bundles" 与 "which are" 之间疑脱 "form" 或 "are"），译文按语义补出。]

\subsection{通过反全纯对合的膜}

通过在紧黎曼面上以及在复李群上同时考虑反全纯对合，可以得到不同类型的膜；这正是 [15, 17] 的观点，我们将在本章余下部分加以描述。仿照 [62] 并回顾第一章中关于实群的 Higgs 丛的分析，对 [89] 中所考虑的 G-Higgs 丛有如下描述：

\begin{prop}{4.4}
设 G 是复半单李群 $G_\C$ 的一个实形式，其实结构为 $\sigma$。则 G-Higgs 丛由 $\cM_{G_\C}$ 中在作用
\[
\Theta_G\colon (P,\Phi)\mapsto(\theta(P),-\theta(\Phi))
\]
下的对合 $\Theta_G$ 的不动点给出，其中 $\theta=\rho\sigma$，$\rho$ 为 $G_\C$ 的紧实形式。
\end{prop}

应当注意，命题 4.4 中 $\cM_{G_\C}$ 里 $\Theta_G$ 的一个不动点给出 $\pi_1(\Sigma)$ 到实形式 G 的一个表示，其中相差 $G_\C$ 中 G 的正规化子的共轭视为等价。这个正规化子可能比李群 G 本身更大，因而 $\cM_G$ 中两个不同的类可以经由一个复映射在 $\cM_{G_\C}$ 中同构。于是，虽然存在从 $\cM_G$ 到 $\cM_{G_\C}$ 中不动点子簇的映射，它却未必是嵌入。秩 2 Higgs 丛情形下这一现象的刻画见 [88, Section 10]，其中可以看到具有不同拓扑不变量的 $\Sl(2,\R)$-Higgs 丛如何与 $\Sl(2,\C)$-Higgs 丛处于同一个连通分支。

% [Tutor Note: p.29 原文 "A description of this phenomena …" 中 "phenomena" 应为 "phenomenon"；同页末段 "appear as the fixed point sets involutions arising from …" 在 "sets" 与 "involutions" 之间疑脱 "of"。两处均照原样理解并译出。]

命题 4.4 的观点（[89] 通篇采用它）被纳入 [15, 17] 所描述的一个更整体的图景：$\cM_{G_\C}$ 中 $(\mathrm{B},\mathrm{A},\mathrm{A})$、$(\mathrm{A},\mathrm{B},\mathrm{A})$ 与 $(\mathrm{A},\mathrm{A},\mathrm{B})$ 型的膜族，作为由复李群上的实形式 $\sigma$ 与黎曼面上的实结构 $f$ 所产生的对合的不动点集而出现。紧连通黎曼面上的实结构\footnote{$\Sigma$ 上的反对合在文献中也称为反共形映射（例如见 [32]）。}是反全纯对合 $f\colon\Sigma\to\Sigma$。

% [原文 p.30]
紧黎曼面上实结构的分类是 Klein 的经典结果：他证明了 $\Sigma$ 上所有这类对合都可由两个整数不变量 $(n,a)$ 来刻画，即：$f$ 的不动点集是嵌入 $\Sigma$ 中的 $n$ 个圆周的不交并；不动点集的补集有一个或两个连通分支，分别对应于 $a=1$ 或 $a=0$。黎曼面 $\Sigma$ 上的实结构 $f$ 在表示 $\pi_1(\Sigma)\to G_\C$ 的模空间、平坦联络的模空间以及 $\Sigma$ 上 $G_\C$-Higgs 丛的模空间上都诱导出对合。

\begin{itemize}[leftmargin=1.4em,itemsep=1pt,topsep=1pt]
\item 由 G 的 Cartan 对合 $\theta$ 得到由
\begin{equation}
i_1(\bar\partial_A,\Phi)=(\theta(\bar\partial_A),-\theta(\Phi)). \tag{4.5}
\end{equation}
固定的 $(\mathrm{B},\mathrm{A},\mathrm{A})$-膜。
\item 由实结构 $f\colon\Sigma\to\Sigma$ 得到由
\begin{equation}
i_2(\bar\partial_A,\Phi)=(f^*(\partial_A),f^*(\Phi^*))=(f^*(\rho(\bar\partial_A)),-f^*(\rho(\Phi))). \tag{4.6}
\end{equation}
固定的 $(\mathrm{A},\mathrm{B},\mathrm{A})$-膜。
\item 最后，考察 $i_3=i_1\circ i_2$，得到由
\[
i_3(\bar\partial_A,\Phi)=(f^*\sigma(\bar\partial_A),f^*\sigma(\Phi))
\]
固定的 $(\mathrm{A},\mathrm{A},\mathrm{B})$-膜。
\end{itemize}

% [Tutor Note: p.30 式 (4.6) 第一项 $f^*(\partial_A)$ 中的 $\partial_A$ 无横杠，而同一等式的第二项写作 $f^*(\rho(\bar\partial_A))$、$i_3$ 式又写作 $f^*\sigma(\bar\partial_A)$；三处对照可知第一项应亦为 $\bar\partial_A$（原文脱横杠），此处照原样排。]

下面我们将从几何与拓扑的角度概述上述这些膜，并说明如何通过 Hitchin 纤维化来研究它们。由 [15, 17] 中上述方法推广而来的结果，如今已在许多其他模空间上获得--例如见 [43, 44, 66, 85] 等。

% [Tutor Note: p.30 原文 "describe how they can be studied though the Hitchin fibration" 中 "though" 应为 "through"（与 p.26 "is though" 同类笔误）；同页脚注 5 的 "as an anti-conformal maps" 中 "an … maps" 亦为语法笔误。]

\para{4.2.1.\ G-Higgs 丛的 (B,A,A)-膜}

正如第一讲中提到的，模空间 $\cM_{G_\C}$ 具有自然的辛结构，我们记之为 $\omega$。此外，按照 [61]，命题 4.4 中的对合 $\Theta_G$（它对应于式 (4.5) 中的 $i_1$）把 $\omega\mapsto-\omega$。因此，在光滑点处不动点集必是拉格朗日的，从而 $\cM_G$ 的预期维数是 $\cM_{G_\C}$ 维数的一半。这些实 G-Higgs 丛的膜可以看成 $\cM_{G_\C}$ 内部、在 Hitchin 底 $\cA_{G_\C}$ 上方那些被保持的纤维中的 $\Theta_G$ 不动点。

借助 Cartan 对经典李代数的分类，我们可以描述在 $G_\C$-Higgs 丛的模空间上可以考虑的各种对合 $i_1$。为此，记 $\mathrm{I}_{p,q}$、$\mathrm{J}_n$ 与 $\mathrm{K}_{p,q}$ 为矩阵
\[
\mathrm{I}_{p,q}=\begin{pmatrix}-\mathrm{I}_p&0\\0&\mathrm{I}_q\end{pmatrix},\qquad
\mathrm{J}_n=\begin{pmatrix}0&\mathrm{I}_n<本地路径>,\qquad
\mathrm{K}_{p,q}=\begin{pmatrix}-\mathrm{I}_p&0&0&0\\0&\mathrm{I}_q&0&0\\0&0&-\mathrm{I}_p&0\\0&0&0&\mathrm{I}_q\end{pmatrix}.
\]

回顾命题 4.4 中的 Cartan 对合 $\Theta_G$ 由 $\theta=\rho\sigma$ 得到，其中 $\rho$ 是 $G_\C$ 的紧实形式，$\sigma$ 是 $G_\C$ 的固定实形式 G 的反全纯对合，我们考虑如下的紧形式与分裂实形式：

% [Tutor Note: p.30 原文 "the cartan involution" 中 "cartan" 首字母小写，应为 "Cartan"。]
